% Appendices B and C: reproduction commands and data formats.
% Evidence: results-tasks.md (F1-F42), results-verify.md (F1-F18),
% results-preprocess.md (P1-P6), and the raw logs. No new measurements.

\section{Reproduction Commands and Key Outputs}
\label{app:reproduce}

All commands run from the project root. The readings come from the rebuild and
rerun of 2026-10-20; raw outputs are archived in \code{reports.zip} under
\code{reports/data/logs/}, and the evidence numbers (F/P) match the text. In
the commands, \code{.../} denotes the same relative path under
\code{reports/data/work/} (\code{tasks/} or \code{verify/}) in the archive; the
complete path appears on the command line of each fact in the results files.

\subsection{Environment, Build, and Tests}

Environment: macOS 27.0.1 (build 26A434), Apple M3, 16\,GiB; CMake 4.4.3
(Unix Makefiles); Apple clang 21.0.0, C++17, \code{-Wall -Wextra}, Release
(\code{-O3 -DNDEBUG}); Catch2 v3 is linked only into the tests;
\code{build/mse} is 329,136 bytes and links only system libraries.

\begin{lstlisting}[style=cli]
cmake -S . -B build -DCMAKE_BUILD_TYPE=Release
cmake --build build -j8
/usr/bin/time -p ctest --test-dir build --output-on-failure
# -> 100% tests passed out of 18
# -> Total Test time (real) =  40.90 sec
\end{lstlisting}

The 18 tests comprise 8 Catch2 unit-test executables (46 \code{TEST\_CASE}s and
85,155 assertions), 2 corpus smoke tests, and 8 independent cross-checks; all
pass in this run.

\subsection{Tasks 1 and 2}

The four word-count configurations were each run three times and the external
wall-clock medians are reported. The commands for stop-word derivation, the
formal index, the three dictionaries, block construction, and the disk sidecar
are as follows.

\begin{lstlisting}[style=cli]
./build/mse count --top 50 --csv .../wordcount-default.csv
./build/mse count --top 50 --stop-words data/stopwords.txt \
    --csv .../wordcount-stop.csv
./build/mse count --top 50 --stem --csv .../wordcount-stem.csv
./build/mse count --top 50 --stop-words data/stopwords.txt --stem \
    --csv .../wordcount-stem-stop.csv
# default: 866626 tokens / 25705 distinct; formal: 556866 / 14937
./build/mse stopwords --derive --min-df-ratio 0.4 --min-entropy-rate 0.93 \
    --out .../derived-df0.40-er0.93.txt --compare data/stopwords.txt
# -> overlap 35 / extra: hath o thee thou / missing: more or than their
./build/mse build --index .../index-formal.bin \
    --stop-words data/stopwords.txt --stem
# -> 195 docs / 14937 terms / 117704 postings / 3685655 bytes / OK
./build/mse build --index .../index-formal-v2.bin \
    --stop-words data/stopwords.txt --stem --compress     # -> 1725694 bytes
./build/mse build --dict splay --index .../index-dict-splay.bin \
    --stop-words data/stopwords.txt --stem                # same for hash/map
./build/mse build-blocks --index .../index-blocks.bin \
    --stop-words data/stopwords.txt --stem
# -> blocks 17 (budget 4096, peak 5129 live)
./build/mse build --index .../formal-v2.bin --dict-file .../formal-v2.bpt \
    --stop-words data/stopwords.txt --stem --compress
./build/mse dump --index .../formal-v2.bin --dict-file .../formal-v2.bpt --verify-all
# -> verified: 14937 terms, 0 mismatches
# -> dict reads: 114 page reads, 29819 cache hits
\end{lstlisting}

\code{--verify-all} is implemented only on the \code{--dict-file} path. A cold
query for \code{banquo} through the sidecar takes 2 page reads and reads 8,308
bytes from the index, 0.225\% of a full load.

\subsection{Task 3, Threshold Experiments, and Independent Cross-Checks}

The query examples and the 201 threshold measurements use the formal index
whose processing flags are \code{stopwords stemmed}. The independent
implementations and the RFC corpus follow; the Python side performs its own
tokenization, index parsing, and token-stream scan.

\begin{lstlisting}[style=cli]
./build/mse query --index .../index-formal.bin --stop-words data/stopwords.txt \
    --word osric --pretty        # -> 1 doc; doc 30 "Hamlet"
./build/mse query --index .../index-formal.bin --stop-words data/stopwords.txt \
    --word love                  # -> 138 documents
./build/mse query --index .../index-formal.bin --stop-words data/stopwords.txt \
    --phrase "king richard" --top-x 3        # -> 3 documents
./build/mse query --index .../index-formal.bin --stop-words data/stopwords.txt \
    --words "king queen" --min-match 2 --top-x 5   # -> 5 documents
./build/mse query --index .../index-formal.bin --stop-words data/stopwords.txt \
    --word love --max-df 0.7     # -> 0 documents; note: df ratio above --max-df
python3 tools/threshold_matrix.py --index .../index-formal.bin \
    --stop-words data/stopwords.txt \
    --csv reports/data/csv/tasks/threshold_matrix.csv --svg .../threshold_matrix.svg
# -> csv: ...threshold_matrix.csv  (201 measurements)

python3 tools/check_wordcount.py --manifest data/manifest.tsv --stem \
    --stop-words data/stopwords.txt > /tmp/wordcount-py.csv
cmp /tmp/wordcount-py.csv .../wordcount-stem-stop.csv       # identical
python3 tools/brute_force.py --manifest data/manifest.tsv \
    --index .../formal-v1.bin --stop-words data/stopwords.txt --stem
# -> 14937 terms, result: OK
python3 tools/check_phrases.py --index .../formal-v1.bin \
    --stop-words data/stopwords.txt --stem --queries 300
# -> 305 sampled / 303 checked / 0 mismatches

./build/mse stats --manifest data2/data/manifest.tsv
./build/mse stopwords --derive --manifest data2/data/manifest.tsv \
    --min-df-ratio 0.4 --min-entropy-rate 0.93 \
    --out .../rfc-stopwords.txt --compare data/stopwords.txt   # -> 58 terms
./build/mse build --manifest data2/data/manifest.tsv \
    --stop-words .../rfc-stopwords.txt --stem --index .../rfc-formal.bin
\end{lstlisting}

\subsection{Bonus: B+ Tree Scaling and the 500,000-Document Corpus}

The B+ tree experiment uses 4\,KiB pages, bulk loading, and 1,000 random
lookups that are all verified.

\begin{lstlisting}[style=cli]
for keys in 100000 1000000 10000000 30000000 50000000; do
  ./build/mse bptree-bench --keys "$keys" --lookups 1000 --out .../bpt-"$keys".bpt
done
# 50000000 keys -> height 4 / 184322 pages / 754987008 bytes
#               -> 1000 verified, 4.00 page reads each
# CSV: python3 .../parse_bptree.py -> reports/data/csv/verify/bptree-bench.csv
\end{lstlisting}

For the 500,000-document synthetic corpus the sequence is generate, count,
block-build v1 and v2, and in-memory build; peak memory is the
\code{maximum resident set size} reported by \code{/usr/bin/time -l}. The
corpus and temporary indexes were deleted after collection.

\begin{lstlisting}[style=cli]
python3 tools/gen_synthetic_corpus.py --root /tmp/mse-scale-corpus-500k \
    --docs 500000 --tokens-per-doc 40 --vocab 5000000
# -> 500000 docs / 20000000 tokens / generate 73.67 s
/usr/bin/time -l ./build/mse stats \
    --manifest /tmp/mse-scale-corpus-500k/data/manifest.tsv   # RSS 463.9 MiB
/usr/bin/time -l ./build/mse build-blocks \
    --manifest /tmp/mse-scale-corpus-500k/data/manifest.tsv \
    --block-terms 4096 --index .../scale-500k-v1.bin
# -> blocks 4855 / terms 4908747 / scan 5.47 s + merge 125.08 s
# -> 427603247 bytes / max RSS 1.01 GiB / real 168.16 s
/usr/bin/time -l ./build/mse build \
    --manifest /tmp/mse-scale-corpus-500k/data/manifest.tsv \
    --index .../scale-500k-mem.bin
# -> 427603247 bytes / max RSS 2.00 GiB / real 76.75 s
# -> cmp with v1: identical (SHA-256 4fcb45...)
\end{lstlisting}

\begin{sloppypar}
The archive \code{reports.zip} organizes the evidence as follows: three result
files (\code{results-tasks.md}, \code{results-verify.md},
\code{results-preprocess.md}); \code{csv/} with the plotting inputs (word
frequency, threshold matrix, B+ tree benchmark); \code{logs/} with 53 raw logs
in three groups (30 + 20 + 3); and \code{work/} with temporary indexes and
reproduction scripts, which are not evidence. Timing convention: two analysts
measured concurrently on the same machine, so wall-clock times may be raised by
parallel load and are meaningful only in order of magnitude. Byte counts,
counts, page counts, SHA-256 values, and cross-check results are deterministic.
\end{sloppypar}

\section{Data Format Reference}
\label{app:format}
\enlargethispage{3\baselineskip}

This appendix summarizes the four data files the engine reads and writes: the
corpus manifest, the text layer, the index (v1 and v2), and the disk B+ tree
sidecar. The format specification follows the implementation and
Sections~\ref{sec:persistence} and~\ref{sec:bptree-design}.

\subsection{Corpus Manifest and Text Layer}

\code{manifest.tsv} is tab-separated with a header row and one document per
line; the \code{text} column is resolved relative to the \emph{corpus root}
(the engine takes \code{dirname(dirname(manifest))}). The optional eighth
column, \code{format}, is \code{marked} (the default, line markers active) or
\code{plain} (the whole line is body text, used by the RFC corpus). The text
layer has one unit per line, and a line-initial marker determines the line
type; headings and stage directions require a match at both the first and the
last character, otherwise the line is body text
(Section~\ref{sec:preprocess-analysis}).

\begin{center}
\footnotesize
\begin{tabularx}{\linewidth}{l l X l}
  \toprule
  Column / marker & Type & Meaning & Default behavior \\
  \midrule
  \code{docid} / \code{kind} / \code{category} & integer / string &
    Document ID, work type, and category & --- \\
  \code{title} & string & Work title, used in output and \code{--pretty} & --- \\
  \code{text} / \code{source} / \code{url} & path / string &
    Text layer, source file path, and source URL & --- \\
  \code{format} & string (optional) & \code{marked} or \code{plain} & --- \\
  \texttt{\#} comment & file header & Excluded from statistics & skipped \\
  \code{[ACT I]} etc. & line starts with \code{[} and ends with \code{]} &
    Heading & not tokenized \\
  \code{(Exit Ghost)} & line starts with \code{(} and ends with \code{)} &
    Stage direction & not tokenized \\
  \code{@BERNARDO} & line starts with \code{@} & Speaker label & not tokenized \\
  Other lines & dialogue, verse, or prose & Body text & tokenized \\
  \bottomrule
\end{tabularx}
\par\smallskip
\footnotesize Note: the three marker classes can be retained independently with
\code{--keep-headings}, \code{--keep-speakers}, and \code{--keep-stage}.
\end{center}

\subsection{Version 1 Index Layout and Version 2 Differences}

All integers are little-endian; a string is a u32 byte length followed by the
raw bytes. The v1 index is laid out sequentially as follows.

\begin{lstlisting}[style=cli]
header : magic "MSE1" | version u32 | flags u32 | docCount u32 |
         termCount u32 | tokenCount u64                 (28 bytes)
doc    : docID u32 | keptTokens u32 | titleLen u32 | title[...]   (12+len)
term   : termLen u32 | term[...] | df u32 | tfTotal u64 |
         blockOffset u64 | blockBytes u32                (28+len)
block  : df x (docID u32, tf u32), then sum(tf) x position u32
         blockBytes = 8*df + 4*sum(tf)
\end{lstlisting}

Bit 0 of \code{flags} records stop-word filtering and bit 1 records stemming;
the query side replays the same normalization pipeline from these flags. The
term table is strictly increasing, document IDs are dense from 1 to \(N\),
document IDs within a posting block strictly increase, and the term frequency
equals the number of positions; \code{validateIndex()} checks these invariants
after loading. The v2 index has the same header, document table, and term table
as v1 except that \code{version}$\,=2$ and each posting block is replaced by
``deltas plus byte-aligned variable-length codes.'' For each posting it writes
the document-ID delta (starting from 0 within the block, $\ge1$), the term
frequency, and the term-frequency position deltas (starting from 0 within the
record, $\ge0$). The first byte's top two bits give the total length and the
payload uses 6/14/22/30 bits (\code{01}/\code{10}/\code{11}/\code{00} for
1/2/3/4 bytes), little-endian.

\subsection{Disk B+ Tree Sidecar}

The sidecar is a 4\,KiB-page file (the page size is configurable). Page 0 is the
header page and each later page is one tree node; integers are little-endian.
Keys are variable-length strings, the fanout is determined by a byte budget, and
leaf pages are linked in order by \code{nextLeaf}.

\begin{lstlisting}[style=cli]
page 0 (header) : magic "BPT1" | version u32 | pageSize u32 |
                  pageCount u32 | rootId u32 | height u32 |
                  leafCount u64 | keyCount u64
page i+1        : leaf u8 | keyCount u16 | nextLeaf u32 |
                  leaf  : value u64[count] | blob[count]
                  inner : child u32[count+1] | blob[count]
blob            : length u8 | bytes
value           : (posting blockOffset << 21) | blockBytes   (21-bit length)
\end{lstlisting}

A query uses the B+ tree to locate the posting block without loading the full
term table. \code{validate()} checks invariants I1--I5. The file is written by
bulk loading, and it is rebuilt whenever the index is rebuilt.
